\section{Deciding a formula and extracting its witnesses}
\label{sec:qe}
Write $K$ for the coded field of the preceding construction. We give the
elimination argument because a decision procedure proved terminating in a
stronger theory would by itself give insufficient information about iPRA.
The geometric atoms are point equality, equidistance and inclusive
betweenness. Translate them respectively into component equality,
equality of squared distances, and
\[
  [C-A,X-A]=0\ \land\ (X-A)\cdot(X-C)\le0.
\]
Section~\ref{sec:synthetic} proves the last equivalence. Relativize numerical
quantifiers to valid codes, use field equality on representatives, replace
each point binder by two scalar binders, and commute the translation with
each intuitionistic connective. This is the interpretation $\varphi^\#$.
Induction on each fixed formula shows that replacement of a parameter by an equal code preserves its interpretation: this holds for atoms, commutes with the connectives, and holds for each admissible assignment under a binder. We next give a quantifier-free ordered-ring formula for each such
formula before interpreting its scalar variables by codes.

\subsection{Counting a sign at real roots}
For a nonconstant squarefree univariate $D$ and any $W$ in $K[T]$, define
\[
   Q_D(W)=\sum_{D(r)=0}\operatorname{sgn} W(r),
\]
with each distinct root counted once. We show how to compute it using
only coefficient arithmetic and sign tests. Reduce $D'W$ modulo $D$;
if the remainder is zero the answer is zero. Otherwise cancel the gcd
of the remainder and $D$, dividing both by the same monic polynomial,
to obtain coprime $B,A$ with $B/A$ the reduced proper fraction. Form
\[
 S_0=A,\quad S_1=B,\qquad S_{i+1}=-\operatorname{rem}(S_{i-1},S_i)
\]
until just before the zero remainder. The last entry is a nonzero
constant. At $+\infty$ record the sign of each leading coefficient;
at $-\infty$ multiply it by $(-1)^{\deg S_i}$. Let $V_+,V_-$ be the
numbers of variations between consecutive signs. Then
\begin{equation}
                  Q_D(W)=V_- - V_+.                 \label{eq:query}
\end{equation}
Indeed, signs are constant between consecutive roots of the product
of the $S_i$ by polynomial intermediate value. Neighboring entries
have no common root, by the Euclidean identities and coprimality.
At a zero of an interior entry the neighboring values are opposite
by the signed remainder identity. Its two adjacent edges have one
variation in total on either side. Several interior entries may vanish
at the same point; their pairs of edges are disjoint. The final entry
has no root. At a root $r$ of the first entry, its root is simple,
and the first edge loses one variation if $B(r)/A'(r)>0$ and gains one
if it is negative. To check this, write $A=(T-r)A_1$, where
$A_1(r)=A'(r)$; sufficiently close to $r$, $A_1$ and $B$ have their
signs at $r$. For a polynomial $g$ with $g(r)\ne0$ this last assertion
has an explicit bound: write $g(r+h)-g(r)=\sum_{j\ge1}c_jh^j$ and take
$|h|<\min(1,|g(r)|/(1+\sum|c_j|))$. Edge changes telescope.
At an uncancelled root of $D$, the coefficient of $1/(T-r)$ in
$D'W/D$ is $W(r)$, as follows from $D=(T-r)D_1$ and $D'(r)=D_1(r)$.
Removing the polynomial part or a common factor preserves this coefficient.
Cancelled roots have $W(r)=0$. This proves~\eqref{eq:query}, including
the case that there are no roots. Signs on the tails have the claimed
values by dominance of the leading term; one may use
$|T|>1+\sum_{j<d}|a_j/a_d|$. For a constant $D\ne0$ set all queries to zero.

For a list $f_1,\ldots,f_m$ and a prescribed sign vector
$\sigma\in\{-1,0,1\}^m$, its count among the roots of $D$ is computed
without comparing unknown roots. Use the three integer polynomials
\[
 I_0(z)=2(1-z^2),\qquad I_{+}(z)=z^2+z,\qquad I_{-}(z)=z^2-z .
\]
On a sign input, each is twice the indicator of its named sign. Expand
$\prod_j I_{\sigma_j}(Z_j)=\sum_e c_e Z^e$. Then the desired count times
$2^m$ is
\begin{equation}
             \sum_e c_e\, Q_D\Big(\prod_j f_j^{e_j}\Big).      \label{eq:counts}
\end{equation}
This follows by expansion at each distinct root and multiplicativity of
sign, with $z^0=1$ also at zero. Positivity of this integer is the
test for occurrence.

\subsection{One scalar quantifier}
Specialize all other variables to any fixed parameter tuple for the
moment. Given the polynomials occurring in a quantifier-free formula
$\theta(T)$, first drop the constants, retaining their signs. Let $P$ be
the monic squarefree part of the product of the rest. If none remain,
one sample $T=0$ suffices. For squarefree extraction over $K$, divide by the gcd with the derivative. At a proposed root $r$ divide repeatedly by $T-r$, with degree fuel, to write the original product as $(T-r)^e U$, $U(r)\ne0$. For $e>0$ its derivative has multiplicity $e-1$ there. The gcd has the smaller multiplicity, by common divisibility and the B\'ezout identity. Thus extraction preserves every root and makes it simple. When a nonconstant polynomial remains, the signs attained anywhere on
the line are precisely those among the following finite tests:
\begin{enumerate}
\item the signs on the left and right unbounded tails;
\item sign vectors occurring at a root of $P$;
\item sign vectors occurring at a root of the squarefree part of $P'$
(omit this test for a constant derivative).
\end{enumerate}
Indeed a sign cannot change inside a root-free interval. For consecutive
roots $a<b$ of $P$, simplicity gives
$\operatorname{sgn}P'(a)=-\operatorname{sgn}P'(b)$: factor $T-a$ at the
left and $T-b$ at the right and use the common interior sign. The
local bound in the preceding paragraph supplies actual one-sided test
points, using also half the gap. Polynomial intermediate value applied
to $P'$ gives a root strictly between $a,b$. Thus the list samples
every bounded open cell as well as all singleton cells. If $P$ has at
most one root, the same proof uses only the singleton and tail cases.
Evaluate the propositional formula on these signs. Existential truth is
equivalent to at least one true evaluation. For witnesses in $K$, use its
sorted root lists directly, including midpoints for bounded intervals and
the extreme roots shifted outwards by one for unbounded intervals
(use zero if there are no roots). An inhabitant and, in the universal
case, a counterexample when it fails, are both primitive recursively
obtainable. This also justifies using Boolean implication for this
specific decidable matrix in intuitionistic logic.

To remove the specialization, run the same coefficient algorithm
symbolically, splitting into three branches at each sign test on the
parameters. In particular test leading terms before dividing, discard
an identically zero specialized polynomial on the appropriate branch,
and truncate to the actual degree. On any branch every scalar coefficient
computed is a rational function $p/q$ in the parameters with $q\ne0$
along that branch. Its sign is the sign of $pq$, since $q^2>0$.
Thus every guard is a polynomial sign formula. The disjunction of the
guards for the accepting leaves is a quantifier-free formula in the
parameters. Keep coefficient pairs formal even on inconsistent branches. A division by $a/b$ is performed symbolically only on the branch guarded by $a\ne0$; inductively all denominators on a realized branch are nonzero. No branch-feasibility oracle is needed.
Symbolic root extraction is unnecessary: counts and tails use only
coefficient arithmetic. This is a specialization-safe elimination even
where a discriminant or several leading coefficients vanish.

Here are bounds on generating the tree. If the original list has $m$
polynomials of degree at most $d$ in $T$, put $s=1+m(d+1)$.
The product has degree below $s$, the squarefree parts and their
derivatives have degree below $s$, and all test products in
\eqref{eq:counts} have degree at most $2md$. A Euclidean gcd or signed
remainder chain uses at most the initial maximum degree plus two
divisions, each a loop of at most that degree plus two cancellations.
Polynomial products use finite convolutions. There are $3^m$ target
patterns and at most $3^m$ monomials per expanded indicator product.
For precision, a finite program with just these bounded loops can be
unrolled by primitive recursion on their displayed bounds. For every
scalar operation the size of a pair of numerator/denominator expression
trees of size at most $b$ is bounded, for instance, by $8(b+1)$ for
each new tree if expressions remain unexpanded. Record a three-way
branch by duplicating the remaining program at most three times per
test. Iteration of these size bounds along the unrolled program, or
construction of its whole ternary tree by recursion on its instruction
budget, is primitive recursive. Expansion and collection of the resulting
polynomials are finite folds on their syntax. Repeating one-variable
elimination from the innermost binder out is iteration of a fixed PR
transformation on coded syntax, hence PR even with formula length an
input. Efficiency is immaterial to the proof.

\begin{theorem}[Elementary equivalence and arithmetical derivability]
\label{thm:transfer}
For each fixed first-order point formula $\varphi(\bar P)$ there is
a PR test $b_\varphi$ such that iPRA proves, on valid parameter codes,
\[
                  \varphi^\#(\bar P)\ \longleftrightarrow\
                  b_\varphi(\bar P)=1 .
\]
The test takes only the values zero and one. For algebraic parameters
the standard real plane agrees with this test. In particular, for a
closed sentence $\sigma$ the following are equivalent: its truth in
the ordinary Euclidean plane; its truth in the coded plane in the
standard naturals; and $\mathrm{iPRA}\vdash\sigma^\#$.
\end{theorem}
\begin{proof}
For the internal assertion use the root and order theorems of the scalar
construction and the elimination calculation above. For each fixed formula,
structural induction here is on the external formula. At an existential
step the accepting branch has a sample witnessing the matrix, and the
rejecting branch excludes every sample and hence every candidate. At a
universal step reject with the constructed counterexample or prove the
matrix for an arbitrary input by its realized sign cell. Equivalently,
one may eliminate $\neg\exists T\,\neg\theta(T)$ after this argument
establishes the equivalence; the witness step matters intuitionistically.
Boolean cases for conjunction, disjunction and implication use the
inductively obtained decisions. No global law of excluded middle enters.
For the finite computations within a step, order and polynomial arithmetic
are PR; variations, sign counts, neighboring-root facts and the telescoping
sum are verified by induction on finite list indices with Boolean
invariants (for the sign-constancy assertion apply IVT to an arbitrary
candidate point with the opposite sign). Parameters, including any
candidate point, stay free through these inductions.

Externally, interpret the same tests over $\mathbb R$. Each algebraic
code denotes its designated real root; the embedding preserves field
operations and signs by the isolation construction. Over the reals the
one-variable argument is identical, using the ordinary polynomial IVT
and the finite sorted real roots of a nonzero polynomial. Thus symbolic
elimination produces the same quantifier-free formula there. All its
atomic signs agree on the given parameters. With no free parameters the
Boolean calculation is a closed PR computation, and iPRA proves its
numeral result by its defining equations. If it evaluates to zero,
iPRA proves the negation of the translation. Conversely, the ordinary
natural numbers satisfy the axioms and rules of iPRA, so a derived
closed translation is true there. This last soundness assertion is made
in the metatheory. The result is a schema of derivations, one per formula;
it asserts no internal soundness or truth predicate for iPRA.
\end{proof}
\noindent
This argument is a direct form of univariate sign determination. For other
constructive implementations see Cohen and Mahboubi~\cite{cm} and for
elementary-recursive complexity bounds see Perrucci and Roy~\cite{pr}.

\paragraph{An elimination one can see in the plane.}
In a unit frame the circles centered at $(0,0)$ and $(1,0)$ with common
squared radius $s\ge0$ meet exactly when
\[
  \exists x\,\exists y\,
  \bigl(x^2+y^2=s\ \land\ (x-1)^2+y^2=s\bigr).
\]
Subtracting the equalities forces $2x=1$. The remaining equation is
$y^2=s-1/4$, so the eliminated condition is the single sign test
$4s-1\ge0$. If it is positive, requiring $y>0$ singles out
$(1/2,\sqrt{s-1/4})$; at zero the only witness is $(1/2,0)$.
At $s=1$ this constructs the upper vertex of the equilateral triangle
on the unit segment. The optional \texttt{witness.py} example computes
both coordinates with exact records. It also finds, for the formula
$\exists x\,\forall y\,(y^2\ne x)$, a negative $x$: witness selection
still works when the matrix after the first binder is itself quantified.
